\documentclass[12pt,a4paper]{article}
\usepackage[utf8]{inputenc}
\usepackage{amsmath,amssymb,physics}
\usepackage{hyperref,geometry,booktabs,xcolor}
\geometry{margin=2.5cm}
\hypersetup{colorlinks=true,linkcolor=blue!70!black,citecolor=blue!70!black}

\newcommand{\gq}{\gamma_q}
\newcommand{\gt}{\gamma_t}
\newcommand{\gph}{\gamma_\phi}
\newcommand{\Tr}{\operatorname{Tr}}
\newcommand{\D}{\Delta}

\title{\textbf{How $T_2$ Enters the Solomon Equations}\\[0.4em]
\large Derivation from the Lindblad Master Equation}
\author{Ege Eroglu \\ TLS-QEC Thesis}
\date{\today}

\begin{document}
\maketitle

%─────────────────────────────────────────────────────────────────────────────
\section{Starting Point: The Full Lindblad Equation}

The density matrix $\rho$ of the qubit+TLS system evolves as:
\begin{equation}
    \dot\rho = -i[H,\rho]
    + \gq(1+n_q)\mathcal{D}[\sigma_-^q]
    + \gq n_q \mathcal{D}[\sigma_+^q]
    + \gt(1+n_t)\mathcal{D}[\sigma_-^t]
    + \gt n_t \mathcal{D}[\sigma_+^t]
    + \frac{\gph^q}{2}\mathcal{D}[\sigma_z^q]
    + \frac{\gph^t}{2}\mathcal{D}[\sigma_z^t]
    \label{eq:lindblad}
\end{equation}
where $\mathcal{D}[L]\rho = L\rho L^\dagger - \tfrac{1}{2}\{L^\dagger L,\rho\}$
and $n_{q,t}$ are thermal photon numbers.

%─────────────────────────────────────────────────────────────────────────────
\section{Equations of Motion for Density Matrix Elements}

Working in the basis $\{|ee\rangle, |eg\rangle, |ge\rangle, |gg\rangle\}$
and focusing on the single-excitation subspace, the relevant elements are:
\begin{itemize}
    \item $\rho_{eg,eg} \equiv P_{eg}$: probability of being in $|eg\rangle$
          (qubit excited, TLS ground)
    \item $\rho_{ge,ge} \equiv P_{ge}$: probability of being in $|ge\rangle$
    \item $\rho_{eg,ge} \equiv \sigma$: the quantum coherence between the
          two single-excitation states
\end{itemize}

From Eq.~\eqref{eq:lindblad}, the equations of motion are:
\begin{align}
    \dot{P}_{eg} &= ig(\sigma - \sigma^*)
        - \Gamma_q P_{eg} + \gt n_t P_{gg} \label{eq:Peg}\\
    \dot{P}_{ge} &= ig(\sigma^* - \sigma)
        - \Gamma_t P_{ge} + \gq n_q P_{gg} \label{eq:Pge}\\
    \dot{\sigma} &= -ig(P_{eg} - P_{ge})
        - i\D\,\sigma
        - \frac{\Gamma_q + \Gamma_t}{2}\,\sigma
        \label{eq:sigma}
\end{align}
where the \emph{total} decay rates including thermal excitation are:
\begin{align}
    \Gamma_q &= \gq(1 + 2n_q) + \gph^q \\
    \Gamma_t &= \gt(1 + 2n_t) + \gph^t
\end{align}

\textbf{Key observation:} The coherence $\sigma$ decays at rate
\begin{equation}
    \boxed{
    \frac{1}{T_2^{eff}} = \frac{\Gamma_q + \Gamma_t}{2}
    = \frac{\gq(1+2n_q) + \gph^q + \gt(1+2n_t) + \gph^t}{2}
    }
    \label{eq:T2eff}
\end{equation}
This is the \emph{effective} $T_2$ of the coherence between $|eg\rangle$
and $|ge\rangle$ — it involves both subsystems and both dephasing rates.

%─────────────────────────────────────────────────────────────────────────────
\section{The Secular Approximation and Solomon Equations}

\subsection{Step 1: Secular approximation}

When the Rabi frequency $\Omega_R = \sqrt{\D^2 + 4g^2}/2$ is large
compared to the decay rates, we can average the coherence equation
\eqref{eq:sigma} over a Rabi period. This is valid when:
\begin{equation}
    \Omega_R \gg \Gamma_q,\,\Gamma_t
\end{equation}

\subsection{Step 2: Adiabatic elimination of coherences}

In the secular approximation, the coherence $\sigma$ reaches its
quasi-steady state on a timescale $T_2^{eff}$, which is fast compared
to the population dynamics. Setting $\dot\sigma = 0$:
\begin{equation}
    \sigma^{ss} = \frac{-ig(P_{eg} - P_{ge})}{i\D + 1/T_2^{eff}}
    \label{eq:sigma_ss}
\end{equation}

\subsection{Step 3: Substituting back}

Substituting $\sigma^{ss}$ into Eqs.~\eqref{eq:Peg}--\eqref{eq:Pge}
and using $P_q = P_{eg}$, $P_t = P_{ge}$:
\begin{align}
    \dot{P}_q &= -\Gamma_q P_q - \Gamma_{exch}(P_q - P_t)
                 + \gt n_t(1 - P_q - P_t) \\
    \dot{P}_t &= -\Gamma_t P_t + \Gamma_{exch}(P_q - P_t)
                 + \gq n_q(1 - P_q - P_t)
\end{align}
where the \textbf{exchange rate} is:
\begin{equation}
    \boxed{
    \Gamma_{exch} = \frac{2g^2 / T_2^{eff}}{\D^2 + (1/T_2^{eff})^2}
    = \frac{2g^2\,\Gamma_{tot}/2}{\D^2 + (\Gamma_{tot}/2)^2}
    }
    \label{eq:Gamma_exch}
\end{equation}
with $\Gamma_{tot} = \Gamma_q + \Gamma_t$.

This is a \textbf{Lorentzian in the detuning} $\D$, with:
\begin{itemize}
    \item \textbf{Width}: $1/T_2^{eff}$ — broader $T_2$ means
          resonance condition is relaxed
    \item \textbf{Peak} (at $\D=0$): $\Gamma_{exch} = 2g^2 T_2^{eff}$
          — longer $T_2$ gives \emph{larger} peak exchange
    \item \textbf{Spectral weight} $\int d\D\,\Gamma_{exch} = \pi g^2$
          — conserved regardless of $T_2$
\end{itemize}

%─────────────────────────────────────────────────────────────────────────────
\section{How $T_2$ Enters Solomon: Physical Summary}

\begin{table}[h!]
\centering
\renewcommand{\arraystretch}{1.5}
\begin{tabular}{lll}
\toprule
Quantity & Formula & Physical role \\
\midrule
$T_2^{eff}$ & $\left(\frac{\Gamma_q+\Gamma_t}{2}\right)^{-1}$ &
    Coherence lifetime; when long, Solomon fails \\
$\gamma_\phi^q$ & $\frac{2}{T_2^q} - \gamma_q$ &
    Pure dephasing; kills coherences without energy loss \\
$\Gamma_{exch}$ & $\frac{2g^2/T_2^{eff}}{\D^2 + (1/T_2^{eff})^2}$ &
    Population exchange rate; Lorentzian in $\D$ \\
Peak $\Gamma_{exch}$ & $2g^2 T_2^{eff}$ &
    Grows with $T_2$ — longer coherence = faster exchange \\
Width & $1/T_2^{eff}$ & Broader at shorter $T_2$ \\
\bottomrule
\end{tabular}
\end{table}

\subsection{Physical interpretation}

$T_2$ enters Solomon through the exchange rate $\Gamma_{exch}$ as the
linewidth of a Lorentzian resonance. The physics is:

\begin{enumerate}
\item \textbf{Short $T_2$ (strong dephasing, $\gamma_\phi$ large):}
    Coherences $\sigma$ die quickly $\to$ secular approximation becomes
    exact $\to$ Solomon is valid. The exchange rate has a broad, low peak.

\item \textbf{Long $T_2$ (weak dephasing, $\gamma_\phi \to 0$):}
    Coherences $\sigma$ survive for many Rabi cycles $\to$
    the adiabatic elimination step \eqref{eq:sigma_ss} fails $\to$
    Solomon is not valid. The exchange rate has a sharp, tall peak.

\item \textbf{The breakdown condition:}
    Solomon fails when $T_2^{eff} \gg T_{Rabi} = \pi/g$, i.e.:
    \begin{equation}
        g \gg \frac{\Gamma_q + \Gamma_t}{2} = \frac{1}{T_2^{eff}}
    \end{equation}
    This is equivalent to $g \cdot T_2^{eff} \gg 1$ — many coherent
    oscillations per decoherence time.
\end{enumerate}

\subsection{Connection to the numerical results}

The threshold $g/\gamma_t \approx 0.07$--$2.7$ found numerically
(scan 08) corresponds precisely to the condition
$g \cdot T_2^{eff} \sim 1$:
\begin{equation}
    g \cdot T_2^{eff}
    = g \cdot \frac{2}{\gq + \gt}
    = \frac{2 \times 0.001}{0.01 + 0.005}
    \approx 0.13 \quad \text{at onset}
\end{equation}
which confirms the theoretical prediction that the breakdown occurs
when approximately one coherent exchange occurs per $T_2$.

%─────────────────────────────────────────────────────────────────────────────
\section{Implications for the T2 Sweep (Script 09)}

The $T_2$ sweep (09\_sweep\_T2.py) tests the prediction:
\begin{equation}
    \text{ISE} \propto (g \cdot T_2^{eff})^n, \qquad n \approx 3.1
\end{equation}
derived from the $g/\gamma_t$ sweep data. As $T_2^{eff}$ decreases
(by increasing $\gamma_\phi$):
\begin{itemize}
    \item Coherences die faster $\to$ ISE decreases
    \item Exchange rate $\Gamma_{exch}$ decreases at resonance
    \item Both models converge to the same incoherent dynamics
    \item Trace distance $D \to 0$
\end{itemize}

The $T_2$ value at which $D$ crosses below 0.05 defines the
\emph{Solomon validity boundary} in the $T_2$ dimension, complementing
the $(g/\Delta, g/\gamma_t)$ phase diagram.

\end{document}
